\section{Wildness of the constructed rectifiers}\label{sec:group-results}

Throughout this section, $K=\Q[a]/(a^2-a+1)$, $b=1-a$, and $R=K[x]$.
We use the identities $ab=1$ and $b^2=-a$.
The tame group $\Tame_3(K)$ is generated by affine and elementary
polynomial automorphisms. A coordinate is called tame if it is a component
of a tame automorphism, and wild otherwise.
Relative tameness over $R$ means tameness in two variables over the ring
$R$; a priori, this is stronger than tameness over the coefficient field
$K$. The deductions from published theorems below are working proofs
within this research branch and are submitted for further verification.

\begin{proposition}[The Shastri rectifier]\label{prop:trefoil-wild}
The coordinate $h_T$ and the constructed rectifier $\Phi_T$ are wild
over $K$.
\end{proposition}
\begin{proof}
In the relative coordinates $(u,v)=(axy+z,-y)$, the base coordinate is
$H=mv+P_x(u)$, where $m=bx^5$, and the coefficient of $u^2$ is $bx$.
The Edo--V\'en\'ereau criterion for length-one coordinates
\cite[Proposition~2]{EV2001} requires every nonlinear coefficient to be
divisible by $m$ for tameness. However, $bx^5\nmid bx$ in $R$.
The change $(y,z)\mapsto(u,v)$ is tame over $R$, so $H$ is wild over $R$.

By the Umirbaev--Yu theorem, a wild $K[x]$-coordinate remains a wild
coordinate of $K[x,y,z]$ in characteristic zero
\cite[Main Theorem]{UY2004}.
The preliminary change $M=(x,y+2,z-5x)$ is tame, so $h_T=H\circ M$
is also wild. A tame automorphism cannot have a wild component;
therefore $\Phi_T$ is wild.
\end{proof}

\subsection{The inner coordinate of a rational-length-two presentation}

For the next lemma, $Y,Z$ denote two variables over $R$.
We call a presentation
\[
 F=p_2Y+Q_2(p_1Z+Q_1(Y))
\]
reduced in the sense of Edo if $p_1$ is monic,
$Q_1\in R[Y]$, $Q_1(0)=Q_2(0)=0$, and
$\gcd(p_1,\text{coefficients of }Q_1)=1$;
here $p_1,p_2\ne0$, $p_2\in K(x)$, and $Q_2\in K(x)[Y]$.
We use the case $\deg Q_1\ge2$, $\deg Q_2\ge2$.

\begin{lemma}[A consequence of Edo's classification]\label{lem:inner-wild}
Suppose that $F\in R[Y,Z]$ is a coordinate with the reduced presentation
above. If $F$ is tame over $R$, then $p_1Z+Q_1(Y)$ is also tame over $R$.
\end{lemma}
\begin{proof}
Here products $\sigma\tau$ are compositions of $R$-algebra
automorphisms: $(\sigma\tau)(f)=\sigma(\tau(f))$.
Edo's Theorem~8(1) gives $F=\sigma\tau(Y)$, where both components of each
of $\sigma,\tau$ have length one; such automorphisms are tame.
As in the proof of Theorem~8(3), normalize
$\sigma(Y)=p_3Z+Q_3(Y)$, where $p_3$ is monic and $Q_3(0)=0$.
This normalization is affine and preserves tameness.
The Russell--Sathaye companion formula has the form
\[
 \sigma(Z)=\frac{\varepsilon}{p_3}
 \bigl(V(p_3Z+Q_3(Y))-Y\bigr),\qquad
 \varepsilon\in R^\times,\quad V\in R[Y].
\]
If $\tau(Y)=p_4Z+Q_4(Y)$, then
\[
 F=-\frac{\varepsilon p_4}{p_3}Y+
 \left(\frac{\varepsilon p_4}{p_3}V+Q_4\right)
 (p_3Z+Q_3(Y)).
\]
After removing the constant term, this is a reduced presentation of the
same polynomial. Its uniqueness, established in Edo's Proposition~2,
gives $p_3=p_1$, $Q_3=Q_1$. Thus the inner polynomial is a component of
the tame automorphism $\sigma$; see also the same normalization in the
proof of Theorem~8(3) \cite{Edo2013}.
\end{proof}

\begin{proposition}[The harmonic model and all constructed $C_N$]
\label{prop:family-wild}
For every $N\ge1$, the coordinate $H_N$ and the constructed rectifier
$\Phi_N$ are wild over $K$. In particular, this holds for the harmonic
model $C_1=\gamma_H$.
\end{proposition}
\begin{proof}
Set
\[
 Q_x(U)=bxU^2+(-1+ax^2)U+g(x).
\]
From $P_x(U)=Q_x(U)-ax^4U$ and $a^2+b=0$, we obtain the identity
\[
 H_N=H(x,p_N+x^Nz,y)
     =-ax^4y+Q_x(ax^{N+1}z+axp_N+y).
\]
Introduce
\[
 \begin{aligned}
 G_N&=x^{N+1}z+by+x\bigl(p_N(x,y)-p_N(x,0)\bigr),\\
 c_N&=axp_N(x,0).
 \end{aligned}
\]
Then
\begin{equation}\label{eq:family-reduced}
 H_N-Q_x(c_N)
 =-ax^4y+\bigl[Q_x(aG_N+c_N)-Q_x(c_N)\bigr].
\end{equation}
This is a reduced presentation as in Lemma~\ref{lem:inner-wild}, with
$p_1=x^{N+1}$, $p_2=-ax^4$, and
$Q_1(y)=by+x(p_N(x,y)-p_N(x,0))$.
The linear coefficient of $Q_1$ is $b$ modulo $x$, so the coprimality
condition holds. The outer polynomial has degree two, and
$\deg_yQ_1\ge2$ because $p_N\equiv2-y^2\pmod x$.

Modulo $x^{N+1}$, the linear coefficient of $Q_1$ is invertible and all
coefficients of $y^j$, $j\ge2$, are nilpotent.
The Russell--Sathaye criterion shows that $G_N$ is an $R$-coordinate.
However,
\[
 [y^2]Q_1=-x\pmod{x^2},\qquad
 x^{N+1}\nmid[y^2]Q_1\quad(N\ge1).
\]
By the length-one criterion, $G_N$ is wild over $R$.
Lemma~\ref{lem:inner-wild} rules out tameness of $H_N-Q_x(c_N)$,
and subtracting a polynomial in $x$ does not change tameness.
The Umirbaev--Yu theorem gives wildness of $H_N$ in three variables.
Since $H_N$ is a component of $\Phi_N$, the latter is also wild.
\end{proof}

This proof applies to every $N$; it is not inferred from testing finitely
many indices. For $N=1$, its inner polynomial is
$G_1=x^2z+by-xy^2$. Computational checks of the normalization and
coefficients verify the calculations but do not replace
Lemma~\ref{lem:inner-wild}.

\section{The standard Cremona subgroup}\label{sec:standard-subgroup}

Fix the notation
\[
 \mathcal C_3(K)=\langle\operatorname{PGL}_4(K),\sigma_3\rangle,
 \qquad
 \sigma_3[X_0:X_1:X_2:X_3]
 =[X_1X_2X_3:X_0X_2X_3:X_0X_1X_3:X_0X_1X_2].
\]
In this section, products of maps act from right to left.
By Blanc--Hed\'en, $\Tame_3(K)\subset\mathcal C_3(K)$.
Their monomial criterion in dimension three states that a monomial
birational map with exponent matrix $M$ belongs to $\mathcal C_3(K)$
if and only if every column sum of $M$ is odd
\cite[Theorem~1.1 and Proposition~6.5]{BH2015}.

\begin{proposition}\label{prop:cremona-membership}
All constructed rectifiers $\Phi_T,\Phi_H,\Phi_N$ belong to
$\mathcal C_3(K)$.
\end{proposition}
\begin{proof}
Consider the subgroup $\mathcal L\subset\Cr_3(K)$ fixing $x$ and
generated by
\[
 \begin{aligned}
 E_y(q)&=(x,y+q(x,z),z),& E_z(q)&=(x,y,z+q(x,y)),\\
 S&=(x,z,y),&B_r&=(x,x^ry,x^{-r}z),
 \end{aligned}
\]
where $q\in K[x,x^{-1},T]$, $r\in\Z$, together with the constant
diagonal maps $(x,cy,dz)$, $c,d\in K^\times$.

All generators belong to $\mathcal C_3(K)$.
For $B_r$, the column sums of the exponent matrix are $(1,1,1)$.
For $D=(x,x^{2s}y,z)$, they are $(1+2s,1,1)$.
Choose $s\ge0$ so that $x^{2s}q$ is a polynomial. Then
\[
 E_y(q)=D^{-1}E_y(x^{2s}q)D.
\]
The middle map is tame. The maps $E_z(q)$ are obtained by conjugation
with $S$; the remaining constant generators are linear.

In the variables $(u,v)$, with $x$ fixed, the base map factors as
\[
 (u,v)\longmapsto\left(P_x(u)+mv,
 \frac{u-F_x(P_x(u)+mv)}m\right)
 =T_FD_mST_{P/m},
\]
where
\[
 \begin{aligned}
 T_{P/m}(u,v)&=(u,v+P_x(u)/m),&D_m(u,v)&=(mu,v/m),\\
 T_F(u,v)&=(u,v-F_x(u)/m),&S(u,v)&=(v,u).
 \end{aligned}
\]
For $m=bx^5$, every factor belongs to $\mathcal L$.
The initial change $(y,z)\mapsto(axy+z,-y)$ is also a product of
generators of $\mathcal L$. Hence $A_0\in\mathcal L$.

For the transfer, set $D_N=(x,y,x^Nz)$ and use the correction polynomial
$\varphi_N(x,T)$ from the construction of the full companion. Then
\[
 \beta_N=SE_z(p_N)D_N,\qquad
 A_N=D_N^{-1}\bigl(E_z(-\varphi_N)A_0SE_z(p_N)\bigr)D_N.
\]
The map $D_N$ normalizes $\mathcal L$, since
\[
 \begin{aligned}
 D_N^{-1}E_y(q)D_N&=E_y(q(x,x^NT)),\\
 D_N^{-1}E_z(q)D_N&=E_z(x^{-N}q),\\
 D_N^{-1}SD_N&=B_NS,
 \end{aligned}
\]
and it commutes with $B_r$ and the constant diagonal maps.
The same formulas with $-N$ give the reverse inclusion.
Therefore $A_N\in\mathcal L\subset\mathcal C_3(K)$.
The outer maps completing the rectifications are tame, which proves
the assertion.
\end{proof}

\begin{remark}
For odd $N$, the map $D_N$ itself fails the monomial criterion:
its first column sum is $1+N$. The proof uses the fact that $D_N$
\emph{normalizes} $\mathcal L$, not membership of $D_N$ in
$\mathcal C_3(K)$.
\end{remark}

\section{Stable tameness and the scope of the conclusion}\label{sec:stable-tameness}

\begin{proposition}\label{prop:stable-tameness}
For each constructed rectifier,
\[
 \Phi_\ast\times\id_{\A^3}\in\Tame_6(K),
 \qquad \ast=T,H,N\ (N\ge1).
\]
\end{proposition}
\begin{proof}
Berson--van den Essen--Wright proved that, in characteristic zero, the
subgroup of $\Aut_K(\A^3)$ generated by automorphisms fixing one
coordinate is contained in $\Tame_6(K)$ after adjoining three variables
\cite[Corollary~4.9]{BEW2012}.
Each of $A_0,A_N$ fixes $x$.
The outer factors of our rectifiers are tame. Elementary automorphisms
fix two coordinates, and the affine group is generated by elementary
linear transformations, scalings of one coordinate, and translations of
one coordinate; each such generator fixes at least one coordinate.
Thus every $\Phi_\ast$ belongs to the stated subgroup.
\end{proof}

Three additional variables suffice; their minimal number is not determined
here. Combining the proofs gives the working conclusion
\[
 \Phi_\ast\in
 \bigl(\mathcal C_3(K)\cap\Aut_K(\A^3)\bigr)
 \setminus\Tame_3(K),\qquad
 \Phi_\ast\times\id_{\A^3}\in\Tame_6(K).
\]
These are properties of the \emph{chosen maps}. They do not rule out
the existence of another tame rectifier of the same curve.

The existence of wild automorphisms in three variables and of wild
coordinates is already established in the literature. Membership of
the wild Nagata automorphism in $\mathcal C_3$ has also been proved
previously \cite[Proposition~6.8]{BH2015}. Accordingly, these results
are not claimed to solve a general tameness problem, the general stable
tameness conjecture, or the problem of generators for the full group
$\Cr_3$.
The content of this application is the simultaneous characterization of
explicit rectifiers of specific curves, together with a proof for every
index of the constructed family.
